\clearpage
\fancyhead[L]{\small VON ZACH. TABULAE ABERRATIONIS ETC.}
\hypertarget{native-494-lower}{}\label{native-494-lower}
\begin{center}
\small Göttingische gelehrte Anzeigen. Stück 2. Seite 18..24. 1808. Januar 2.\\[3pt]
\rule{.72\linewidth}{.3pt}
\end{center}

\emph{Tabulae speciales aberrationis et nutationis in ascensionem rectam et in declinationem ad supputandas stellarum fixarum positiones, sive apparentes, sive veras, una cum insigniorum 494 stellarum zodiacalium catalogo novo in specula astronomica Ernestina ad initium anni \textup{1800} constructo, cum aliis tabulis eo spectantibus, auctore} \textsc{Franc. Lib. Bar. de Zach}. Vol. I. 1806. 208 und \textsc{clvi} S. Vol. II. 1807. 508 S. gr. Octav.

Die Erscheinung dieses wichtigen, die kostbaren Früchte vieljähriger Arbeit enthaltenden, Werks wird allen practischen Astronomen, die dasselbe schon seit mehreren Jahren erwarteten, höchst willkommen sein. Sein Hauptgegenstand bezieht sich auf die kleinen, von Präcession, Aberration und Nu%
\SourceLeaf{495}%
tation herrührenden, Veränderungen der Lage der Fixsterne. Diese Lage ist bekanntlich die Basis der meisten astronomischen Beobachtungen: die Berechnung von jenen kleinen Aenderungen gehört daher zu den täglich wiederkehrenden Beschäftigungen des practischen Astronomen. Um diese hat sich also der berühmte Verfasser durch ein Werk sehr verdient gemacht, das die geschmeidigsten und brauchbarsten Formeln für jene Veränderungen aufstellt, durch die sorgfältigsten Erörterungen nach den neuesten und bewährtesten Untersuchungen und Beobachtungen die zuverlässigsten numerischen Werthe der dabei zum Grunde liegenden Grössen entwickelt, und durch zweckmässige und bequeme Hülfstafeln die Arbeit so viel als möglich abkürzt, also, mit Einem Worte, alles erschöpft, was über diesen Gegenstand nur gewünscht werden kann. So sehr schätzbar alles dieses ist, so möchte Rec. doch einen andern beträchtlichen Theil des Werks, der auf dem Titel nicht einmal erwähnt ist, noch höher anschlagen, nemlich das Verzeichniss von 1830 Zodiacal-Sternen, ein monumentum aere perennius von zwei Sternwarten, welche die Zierde von Deutschland waren, und jetzt, leider! beide unbenutzt stehen. Einem Werke von dieser Wichtigkeit müssen wir eine ausführlichere Anzeige widmen.

Die mehr als die Hälfte des ersten Bandes einnehmende Einleitung macht mit der Präcession den Anfang. Nach einer vorausgeschickten kurzen Erzählung der ersten Entdeckung derselben, und der theoretischen Arbeiten von \textsc{Newton}, \textsc{d'Alembert}, \textsc{Euler}, \textsc{Lagrange} und \textsc{Laplace} wird eine Tafel für die Ungleichheit der Präcession in den verschiedenen Jahrhunderten gegeben; hierauf Formeln, und durch einen dem Verf. eigenthümlichen Kunstgriff bequem eingerichtete Hülfstafeln für die Säcularänderung der Breite, und denjenigen Theil der Säcular-Aenderung der Länge, der für verschiedene Sterne verschieden ist. Hierauf folgen die Formeln für die Präcession in Rectascension und Declination, nebst einer Tafel für die Werthe der Constanten vom J. 1450 bis 1950. Die Fundamentalwerthe hat der Verfasser durch eigne sehr sorgfältige Vergleichungen ausgemittelt, wozu er mit Recht blos die neuern Beobachtungen gebraucht hat. Denn \textsc{Hipparch}'s und \textsc{Tycho}'s Beobachtungen können, der grössern Zwischenzeit ungeachtet, doch ihrer viel geringern Vollkommenheit wegen, bei weitem keine so zuverlässige Resultate geben, als die mit den besten Mauerquadranten angestellten von \textsc{Mayer} und \textsc{Bradley}. Auch selbst \textsc{Flamsteed}'s Beobachtungen stehen diesen noch zu sehr nach, und würden auf alle Fälle, wenn etwas Brauchbares aus ihnen gefolgert werden sollte, nothwendig einer ganz neuen Reduction bedürfen. Dass indess die etwa noch vorhandene kleine Ungewissheit in der Abnahme der Schiefe der Ekliptik und in der eignen Bewegung der Fixsterne bei entferntern Beobachtungen mehr betrage, als bei nähern, würde Rec. unter den Gründen, warum die \textsc{Flamsteed}'schen Beobachtungen ausgeschlossen wurden, nicht angeführt haben. Aus der Vergleichung von 175 Declinationen von \textsc{Bradley} mit denen von \textsc{Barry} leitet der Verf. die ganze jährliche Präcession $50''1670$ ab; aus Vergleichung der Declinationen von \textsc{Mayer} und \textsc{Barry} $49''9140$; ferner gaben 350 Rectascensionen von \textsc{Mayer}, verglichen mit des Verf. eigner Bestimmung, $50''0718$, und endlich 222 Rectascensionen von \textsc{Bradley} mit denen des Verf. $50''0631$, also im Mittel aus allen Bestimmungen $50''0540$, welcher Werth bei allen Formeln und Tafeln des Werks zum Grunde gelegt ist. Wir berühren nur die hierauf folgenden Untersuchungen über die eignen Bewegungen der Sterne, wo man unter andern mit Vergnügen eine ganz neue Reduction aller Beobachtungen der Declination des Polarsternes von \textsc{Flammsteed} finden wird, woraus sich blos eine Säcular-Zunahme von $3''$ ergibt. Zur Berechnung der Declination aus Länge und Breite gebraucht hier der Verf. Tafeln mit 10 Decimalstellen: dies ist freilich nothwendig, wenn man die Declination durch den Sinus bestimmt, allein dieser Inconvenienz kann man ganz bequem ausweichen, und mit den gewöhnlichen Tafeln eben so weit reichen. Eine Tafel für die mittlere Rectascension und Declination des Polarsterns für alle einzelnen Jahre von 1790 bis 1820 macht den Beschluss dieser reichhaltigen Abhandlung.

\begingroup\rightskip=0pt plus 1em
Zur Bestimmung der Constante der \emph{Aberration} hat der Verf. die sämmtlichen \textsc{Bradley}'schen Ori%
\SourceLeaf{496}%
ginalbeobachtungen aufs neue in grösster Schärfe discutirt: er findet dieselbe $20''232$, welches von dem aus den Verfinsterungen der Jupiterstrabanten von \textsc{Delambre} bestimmten Resultate, nemlich $20''255$, nur ganz unbedeutend abweicht. Letztern Werth hat der Verfasser bei einer am Ende des zweiten Bandes befindlichen Generaltafel für die Aberration zum Grunde gelegt. Diese Tafel nimmt nur zwei Seiten ein; ihr Gebrauch erfordert aber Multiplication durch einen Sinus und eine Secante. Auch dieser kleinen Unbequemlichkeit hat der Verf. durch einen eignen Kunstgriff abgeholfen, so dass der Gebrauch einer andern, in die Einleitung eingerückten, Generaltafel blos Additionen erfordert. Hierauf folgen noch Formeln und Tafeln für die Aberration, mit Rücksicht auf die elliptische Bewegung der Erde, und für die Aberration, wegen der täglichen Bewegung; endlich Formeln und Tafeln für die Veränderung der Aberration bei $1^\circ$ Aenderung in Rectascension und Declination des Sterns, wonach sich die Dauur der Brauchbarkeit von Specialtafeln schätzen lässt.\par\endgroup

Auch für die \emph{Nutation} hat der Verf. ausser den Formeln, zwei Generaltafeln geliefert. Die eine am Ende des zweiten Bandes, wo die halbe grosse und kleine Axe der Nutationsellipse nach \textsc{Laplace} zu $10''056$ und $7''486$ vorausgesetzt sind, erfordert Multiplication durch eine Tangente: diese Unbequemlichkeit wird in der andern der Einleitung einverleibten vermieden, wo aber Hr. \textsc{von Zach} jene halben Axen zu $9''648$ und $7''182$ angenommen hat, weil er die von \textsc{Laplace} zum Grunde gelegte Mondmasse etwas vermindern zu müssen glaubte. Ferner findet man in diesem Abschnitte Vorschriften zur Berechnung des von der Sonnenlänge abhängigen Theils der Nutation, und Formeln und Tafeln für die Aenderung der Nutation bei veränderter Lage des Sterns.

Den übrigen Theil der Einleitung füllen Vorschriften zur richtigen Stellung des Passageinstruments; Erklärung der in diesem Bande vorkommenden Verzeichnisse und Tafeln; Untersuchungen und Tafeln in Beziehung auf die Aberration der Planeten, ihre Parallaxe und Durchmesser. Die Behauptung S. 203, dass die (unsrer Meinung nach allen andern vorzuziehende) Art, über die Aberration der Planeten Rechnung zu führen, indem man nemlich nur die Zeit ändert, --- blos den von der Bewegung des Planeten abhängigen Theil der Aberration gebe, können wir nicht beipflichten, wohlverstanden, dass auch der dabei anzuwendende Ort der Erde dem geänderten Zeitmoment entsprechen muss.

\begingroup\rightskip=0pt plus 1em
Uebrigens sind alle in dieser gehaltreichen Einleitung gegebenen Vorschriften durch zahlreiche wohlgewählte Beispiele so erläutert, dass auch der Ungeübteste bei ihrer Anwendung keinen Anstoss finden wird, und überall die Schriften, wo man sich über die abgehandelten Materien weiter belehren kann, nachgewiesen.\par\endgroup

Von den nun folgenden Tafeln selbst zeigen wir nur summarisch den Inhalt an. Die geraden Aufsteigungen und Abweichungen der 36 vornehmsten Sterne nach \textsc{Maskelyne}'s neuester Bestimmung, nebst speciellen Aberrations- und Nutationstafeln für jeden derselben; ähnliche Tafeln für den Polarstern für 1790, 1800, 1810 und 1820; ein Verzeichniss von Fixsternen von fast gleichen Rectascensionen und fast gleichen aber entgegengesetzten Declinationen zum Behuf der Berichtigung der Stellung des Passageinstruments. Hierauf folgt der kostbarste Theil des ganzen Werks, nemlich die auf der Seeberger Sternwarte bestimmten Rectascensionen von 1830 Zodiacalsternen für das Jahr 1800, nebst den Unterschieden von \textsc{Piazzi}'s Catalog; diese Unterschiede sind meistens nur klein, und werden in der Regel noch viel kleiner, wenn man \textsc{Piazzi}'s Rectascensionen die Verbesserung von $3''8$ hinzufügt; immer aber wird man berechtigt sein, den mit dem prächtigen 8füssigen Seeberger Mittagsfernrohr gemachten Bestimmungen den Vorzug zu geben. Die Declinationen sind hier, blos um die Sterne zu designiren, nur in Minuten angegeben; bei einigen Sternen fehlen sie aber ganz; es wäre zu wünschen, dass Jemand die Mühe übernähme, diese Lücke auszufüllen, wobei meistens die Histoire céleste schon ausreichen würde, insofern nicht zufällig mehrere Sterne im Zodiacus auf eine Zeitsecunde in der geraden Auf%
\SourceLeaf{497}%
steigung übereinstimmen; viele von diesen Sternen scheinen sogar in dem gleichfolgenden \textsc{Barry}'schen Declinationsverzeichnisse vorzukommen (z.\,B. in der Jungfrau 921 und 963). Um nun auch eine hinlängliche Anzahl von scharf bestimmten Declinationen zu geben, hat Hr. \textsc{von Zach} ein zweites blos diese nebst der beiläufigen Rectascension enthaltendes Verzeichniss von etwa 1200 Zodiacalsternen beigefügt, welches das Resultat von den auf der Mannheimer Sternwarte von \textsc{Barry} und \textsc{Henry} mit einem 8füssigen \textsc{Bird}'schen Mauerquadranten gemachten Beobachtungen darstellt. Die zugleich beigefügten Unterschiede von \textsc{Piazzi}'s Cataloge zeigen, dass jenes Verzeichniss diesem wohl zur Seite gesetzt zu werden verdient. Auch in diesem Verzeichnisse fehlen bei verschiedenen Sternen die Rectascensionen ganz; dieser Mangel ist hier aber weniger von Bedeutung, da in den meisten Fällen schon die Ordnung zur Erkennung des Sterns dienen kann.

Der noch übrige Theil des ersten Bandes enthält eine allgemeine Tafel für die Präcession in gerader Aufsteigung; Tafeln für die mittlern Refractionen und ihre Verbesserung wegen des Barometer- und Thermometerstandes, nach denselben Grundsätzen, wie die Tafeln bei \textsc{Delambre}'s neuen Sonnentafeln, nur auf die bisher allgemein üblichen Maasse reducirt; Sonnenparallaxe; Verwandlung der Sternzeit in mittlere und wahre Sonnenzeit, und einige andere kleinere Tafeln.

Den \emph{zweiten Band} füllen fast ganz die speciellen Aberrations- und Nutationstafeln für 494 Zodiacalsterne, nebst ihren Positionen für 1800 nach dem Verf. u. a. Astronomen. Wie höchst schätzbar dieses Hülfsmittel den practischen Astronomen sei, haben wir schon oben erwähnt; es bleibt uns also blos zu bemerken übrig, dass dieser zweite Band schon früher gedruckt ist, als der erste, und daher in verschiedenen Punkten von diesem abweicht. Die von dem Verf. herrührenden Sternpositionen müssen daher sämmtlich gegen die in dem grossen Catalog des Bandes befindlichen vertauscht werden, welche nach \textsc{Maskelyne}'s Verbesserung seines Fundamentalverzeichnisses berichtigt sind. Sollen ferner die Präcessionen, Aberrationen und Nutationen mit den neuesten Bestimmungen im ersten Bande ganz übereinstimmend gemacht werden, so müssen die ersten um $\frac{1}{124}$ vermindert, die zweiten um $\frac{1}{78}$ vermehrt, die dritten um $\frac{1}{92}$ vermehrt werden.

Den Beschluss des zweiten Bandes machen die schon oben erwähnten Generaltafeln für Aberration und Nutation, und eine Tafel für die Länge des aufsteigenden Knoten der Mondbahn nach den \textsc{Bürg}'schen Mondtafeln.
